\subsection{Subsets and Power Sets}

\begin{definition} \label{def:proper-subsets-singletons-and-nonempty} \hlt{Basic Subset Terminology}\\
Let $A,B$ be sets. Write $A\subseteq B$; $A=B$ iff $A\subseteq B$ and $B\subseteq A$.
\begin{enumerate}[label=(\roman*)]
\setlength{\itemsep}{0pt}
\item $A$ is a \hlt{proper subset} of $B$, written $A\subsetneq B$, iff $A\subseteq B$ and $A\neq B$.
\item $A$ is a \hlt{singleton} iff $\exists x$ such that $A=\{x\}$.
\item $A$ is \hlt{nonempty} iff $A\neq\emptyset$.
\end{enumerate}
\end{definition}

\begin{method} \label{meth:proof-by-containment} \hlt{Proof by Two Containments}\\
To prove $A=B$, prove both inclusions.
\begin{enumerate}[label=(\roman*)]
\setlength{\itemsep}{0pt}
\item To prove $A\subseteq B$, let $x\in A$ be arbitrary and use the definitions and hypotheses to show $x\in B$.
\item To prove $B\subseteq A$, let $x\in B$ be arbitrary and use the definitions and hypotheses to show $x\in A$.
\item Once both containments are proved, conclude $A=B$ by extensionality of sets.
\end{enumerate}
If both sides are described by membership conditions, this is a chain, $x\in A \Longleftrightarrow \cdots \Longleftrightarrow x\in B$.
\end{method}

\begin{definition} \label{def:power-set} \hlt{Power Set}\\
For a set $X$, the \hlt{power set} is $\mathcal{P}(X)\defeq\{S:S\subseteq X\}$.
\end{definition}

\begin{example} \label{ex:small-power-sets} \hlt{Small Power Sets}
\begin{enumerate}[label=(\roman*)]
\setlength{\itemsep}{0pt}
\item If $X=\emptyset$, then $\mathcal{P}(X)=\{\emptyset\}$.
\item If $X=\{a,b\}$ with $a\neq b$, then $\mathcal{P}(X)=\{\emptyset,\{a\},\{b\},X\}$.
\end{enumerate}
Hence $S\in\mathcal{P}(X)$ iff $S\subseteq X$; membership in $\mathcal{P}(X)$ is containment in $X$.
\end{example}
